\meetingheader
    {Mar 06, 2026}
    {}
    {\meetinglinks{
        \meetinglink{OpenAI Harmony Docs}{https://developers.openai.com/cookbook/articles/openai-harmony/}
        \meetinglink{Hamel Husain -- Tool Injection Format}{https://hamel.dev/notes/llm/openai/func\_template.html}
        \meetinglink{Anthropic Skill Creator}{https://github.com/anthropics/skills/tree/main/skills/skill-creator}
        \meetinglink{ACE Framework}{https://github.com/ace-agent/ace}
        \meetinglink{MCP-Bench}{https://github.com/Accenture/mcp-bench}
    }}
    {\begin{itemize}
        \item Debugging agent pipeline; adapting VLMEvalKit to use ARC API judge
        \item Context evaluation for agentic systems -- literature review
        \item How tool descriptions are injected into the system prompt (OpenAI format confirmed)
        \item Gap: no framework for optimizing tool descriptions for agent performance
        \item Potential research direction: context valuation for tool descriptions (Shapley-style)
    \end{itemize}}
    {}


%% ============================================================
%% SECTION 1: HOW TOOL DESCRIPTIONS ENTER THE CONTEXT WINDOW
%% ============================================================

\section{How Tool Descriptions Enter the Context Window}

Tool descriptions are \textbf{injected directly into the system/developer message} as input tokens on every API call.
OpenAI's Harmony docs\footnote{\url{https://developers.openai.com/cookbook/articles/openai-harmony/}} confirm the exact format, and jailbreak extractions\footnote{\url{https://hamel.dev/notes/llm/openai/func\_template.html}} confirm it matches production GPT-4/4.1.

\subsection{API request vs.\ model input}

\paragraph{What the developer sends (API request):}

\begin{verbatim}
{
  "tools": [{
    "type": "function",
    "function": {
      "name": "get_weather",
      "description": "Gets the weather for a location.",
      "parameters": {
        "type": "object",
        "properties": {
          "location": { "type": "string" },
          "unit": { "type": "string", "enum": ["celsius", "fahrenheit"] }
        },
        "required": ["location"]
      }
    }
  }]
}
\end{verbatim}

\paragraph{What the model actually sees (injected into developer message):}

\begin{verbatim}
<|start|>developer<|message|>
# Instructions
Use a friendly tone.

# Tools

## functions

namespace functions {

// Gets the weather for a location.
type get_weather = (_: {
  // The city and state, e.g. San Francisco, CA
  location: string,
  format?: "celsius" | "fahrenheit",
}) => any;

} // namespace functions
<|end|>
\end{verbatim}

\noindent Key observations:
\begin{itemize}
    \item Tool descriptions become \textbf{TypeScript-like type definitions} inside a \texttt{namespace functions \{\}} block
    \item The \texttt{description} field becomes a \texttt{//} comment above the function signature
    \item Parameter descriptions also become inline \texttt{//} comments
    \item Every API call \textbf{re-sends this entire block} as input tokens
    \item Verbose descriptions on many tools $\times$ many agent steps $=$ massive token overhead
\end{itemize}

\subsection{Implication for DCA}

This maps directly onto DCA's tool-use failures.
With 15+ tools, each tool description is re-injected on every turn of the agent loop.
The quality of these descriptions determines whether the LLM can correctly select and parameterize tools.
Our dependency injection (Section~\ref{sec:injection}, Feb~27) already reduces the \emph{parameter schema} the LLM sees -- but the \emph{description text} is equally important and currently unoptimized.


%% ============================================================
%% SECTION 2: LITERATURE REVIEW -- CONTEXT EVALUATION
%% ============================================================

\newpage
\section{Literature Review: Context Evaluation for Agentic Systems}

Emerging literature shows that context files (AGENTS.md, CLAUDE.md, system prompts, tool descriptions) affect agent performance -- but untested context often \emph{hurts} rather than helps.

\subsection{Anthropic's approach (Skill-Creator)}

Anthropic ships a skill-creator tool (\texttt{anthropics/skills} on GitHub) that solves context evaluation for Claude Code skills:
\begin{itemize}
    \item Counterfactual A/B testing: runs tasks with and without context in parallel
    \item Assertion-based grading: objectively verifiable checks on outputs
    \item Benchmark aggregation: pass rates, timing, token usage, deltas
    \item Human-in-the-loop: browser eval viewer for qualitative feedback
    \item Optional blind comparison via comparator agent
\end{itemize}
This is essentially \textbf{unit testing for context} -- but only works for Claude Code skills.

\subsection{Relevant literature}

\begin{table}[H]
\centering
\small
\caption{Context evaluation and tool description literature.}
\label{tab:context-lit}
\begin{tabular}{p{4.5cm}p{7.5cm}l}
\toprule
\textbf{Paper} & \textbf{Key Finding} & \textbf{Ref} \\
\midrule
Evaluating AGENTS.md (ETH Zurich, Feb 2026) & Context files reduce task success rates, increase cost 20\%+ & \cite{evaluatingAgentsMd2026} \\
\addlinespace
Impact of AGENTS.md (JAWs/ICSE 2026) & Replication with Codex -- same negative finding & \cite{impactAgentsMd2026} \\
\addlinespace
SWE Context Bench (Feb 2026) & Filtered context helps, unfiltered hurts -- retrieval quality matters & \cite{sweContextBench2026} \\
\addlinespace
Agent READMEs (Nov 2025) & 2,303 files analyzed -- median readability = legal contract level & \cite{agentReadmes2025} \\
\addlinespace
ACE: Agentic Context Engineering (Oct 2025) & Evolving context playbooks: +10.6\% on agents, +8.6\% on finance & \cite{ace2025} \\
\addlinespace
MCP-Bench (Accenture, 2025) & Benchmarks tool-level schema understanding and usage & \cite{mcpBench2025} \\
\addlinespace
MCPAgentBench (Dec 2025) & Real-world MCP tool use benchmark with difficulty awareness & \cite{mcpAgentBench2025} \\
\addlinespace
Smelly Tool Descriptions (Feb 2026) & 97.1\% of MCP tool descriptions have smells; Purpose + Guidelines is highest-leverage fix; component effectiveness is domain-dependent & \cite{smellyToolDesc2026} \\
\addlinespace
A Survey of Context Engineering (Jul 2025) & Comprehensive taxonomy from 1,400+ papers & \cite{contextEngineeringSurvey2025} \\
\bottomrule
\end{tabular}
\end{table}

\subsection{The gap}

Nobody has a framework for: given a tool (MCP server, AutoGen tool, function schema), \textbf{what is the optimal description so an agent can use it correctly?}

This is a data-centric problem. A data-centric approach would:
\begin{enumerate}
    \item Decompose tool descriptions into atomic instruction components
    \item Measure the marginal contribution of each component (Shapley-style)
    \item Automatically curate/prune descriptions based on measured utility
\end{enumerate}


%% ============================================================
%% KEY FINDINGS
%% ============================================================

\keyfindings{
    \item \textbf{Tool descriptions are injected into the system prompt as TypeScript-like type definitions on every API call.}
        Confirmed via OpenAI Harmony docs and jailbreak extractions matching production GPT-4/4.1 format.
    \item \textbf{No existing framework optimizes tool descriptions for agent performance.}
        Context evaluation work (AGENTS.md papers, SWE Context Bench) focuses on repository-level context, not tool-level descriptions.
    \item \textbf{97.1\% of MCP tool descriptions have ``smells,'' and the optimal fix is domain-dependent.}
        Purpose + Guidelines dominates in Finance; Purpose + Examples in 3D Design.
        Implies DCA tool descriptions should adapt to the data domain being curated.
    \item \textbf{DCA's dependency injection already optimizes the parameter schema} (Feb~27).
        The next optimization target is the description text itself.
}
